\chapter*{پیش‌گفتار}
\addcontentsline{toc}{chapter}{پیش‌گفتار}

این کتابچه، بخشی مستقل از \BookletSubtitle\ است --- کتابی درباره‌ی
آرایه‌ها که از یک اختلاف ولتاژ روی یک ترانزیستور شروع می‌شود و تا
هسته‌های تنسور و بردارهای حالت کوانتومی پیش می‌رود. اما پیش از رسیدن
به هر کدام از این‌ها، باید اول یک سؤال پاسخ داده شود: داده واقعاً
چیست، و چطور چیزی انتزاعی سرانجام به‌وسیله‌ی چیزی فیزیکی حمل می‌شود؟

این همان سؤالی است که این کتابچه در سه فصل دنبال می‌کند. فصل اول
می‌پرسد داده چیست، با کمک بخشی کوتاه از \lr{Eighteenth Brumaire}
اثر مارکس نشان می‌دهد که یک رشته‌ی ثابت از کاراکترها برای یک کودک،
یک مورخ و یک اقتصاددان سیاسی داده‌ای کاملاً متفاوت است، و با گذر از
هرم \lr{DIKW} (داده، اطلاعات، دانش، خرد) به پرسشی سخت‌تر می‌رسد: یک
کلمه، یک عدد یا یک نشانه اصلاً چطور جای چیز دیگری را می‌گیرد. فصل دوم
مستقیم سراغ همین پرسش می‌رود و می‌پرسد بازنمایی از کجا می‌آید، با
تکیه بر هیپوکامپ، حافظه‌ی رویدادی و چهار معیار برای چیزی که بازنمایی
نامیده می‌شود، و با پرسشی دیگر پایان می‌یابد: آیا بازنمایی می‌تواند
به‌صورت یک سامانه‌ی فیزیکی هم وجود داشته باشد، نه فقط ذهنی. فصل سوم
دقیقاً از همان‌جا ادامه می‌دهد و این پرسش را تا انتها دنبال می‌کند ---
از یک جانشین انتزاعی برای چیزی دیگر، تا یک سیگنال فیزیکی که واقعاً
می‌تواند آن را حمل کند.

\vfill

\begin{flushleft}
\textit{\BookletAuthor} \\
\textit{\BookletYear}
\end{flushleft}
